\documentclass[12pt,letterpaper]{article}
\usepackage[utf8]{inputenc}
\usepackage[spanish]{babel}
\usepackage{amsmath, amssymb, amsfonts}
\usepackage{graphicx}
\usepackage{hyperref}
\usepackage{enumitem}
\usepackage[margin=2.5cm]{geometry}

\title{Sesión 2: IEEE 754 y errores numéricos}
\author{Análisis Numérico \\ \small{Universidad de la Nopalera}}
\date{}

\begin{document}

\maketitle

\begin{center}
\textbf{Objetivo:} Entender cómo la compu guarda los números y por qué a veces se pone loca.
\end{center}

\section*{1. ¿Cómo se guarda un número real en la computadora?}

En la vida real usamos decimales, pero la computadora solo entiende binario. Para guardar números con punto decimal, usa el estándar \textbf{IEEE 754}.

Un número en este estándar se guarda como:

\[
\text{Número} = (-1)^s \times (1 + m) \times 2^{e - \text{bias}}
\]

Donde:
\begin{itemize}
    \item $s$ es el bit de signo (0 para positivo, 1 para negativo).
    \item $m$ es la mantisa (los decimales en binario).
    \item $e$ es el exponente (un entero guardado con un sesgo o \textit{bias}).
\end{itemize}

En \textbf{doble precisión} (el que usan casi todos los programas), los tamaños son:

\begin{center}
\begin{tabular}{|c|c|}
\hline
\textbf{Componente} & \textbf{Bits} \\
\hline
Signo & 1 \\
Exponente & 11 \\
Mantisa & 52 \\
\hline
\end{tabular}
\end{center}

\subsection*{Ejemplo: $0.3$ en binario}
Ya vimos que:
\[
0.3_{10} = 0.01001100110011\ldots_2
\]
En notación científica binaria:
\[
0.3_{10} = 1.001100110011\ldots_2 \times 2^{-2}
\]
La compu guarda solo los primeros 52 bits de la mantisa. El resto se pierde. ¡Ese es el origen del error!

\section*{2. El problema de $0.1 + 0.2$}

Tanto $0.1$ como $0.2$ tienen representación binaria \textbf{infinita y periódica}. La compu los redondea al guardarlos.

Cuando los suma, el error de cada uno se acumula y el resultado no es exactamente $0.3$.

En Julia (y en casi todos los lenguajes):
\begin{verbatim}
julia> 0.1 + 0.2 == 0.3
false
\end{verbatim}

\section*{3. Errores que truenan métodos numéricos}

\subsection*{3.1 Resta de números casi iguales (cancelación)}
Si calculas $\sqrt{x+1} - \sqrt{x}$ para $x$ muy grande, los dos números son casi iguales y la resta pierde precisión.

\textbf{Solución:} Usar la identidad equivalente:
\[
\sqrt{x+1} - \sqrt{x} = \frac{1}{\sqrt{x+1} + \sqrt{x}}
\]

\subsection*{3.2 Evaluación de polinomios}
El polinomio $P(x) = (x-1)^{10}$ se puede abrir como:
\[
x^{10} - 10x^9 + 45x^8 - \cdots + 1
\]
Si evalúas en $x = 0.999$, la forma desarrollada da errores enormes. La forma factorizada es mucho más estable.

\subsection*{3.3 Orden de la suma}
La suma en punto flotante \textbf{no es asociativa}. Sumar primero los números más pequeños evita que se pierdan al sumarlos con los grandes.

\section*{4. Ejercicios para la sesión}

\begin{enumerate}
    \item En Julia, calculen:
    \begin{verbatim}
    0.1 + 0.2 == 0.3
    0.1 + 0.2 - 0.3
    \end{verbatim}
    Expliquen el resultado.

    \item Evalúen $(1.001 - 1)^6$ usando la forma factorizada y la forma desarrollada. Comparen.

    \item Sumen $\sum_{n=1}^{1000} \frac{1}{n^2}$ en dos órdenes diferentes:
    \begin{itemize}
        \item De $n=1$ a $1000$.
        \item De $n=1000$ a $1$.
    \end{itemize}
    Comparen con $\pi^2/6 \approx 1.644934$.
\end{enumerate}

\section*{5. Frase para cerrar}

\begin{center}
\fbox{\begin{minipage}{0.9\textwidth}
\centering
\textit{“En análisis numérico no se trata de que la compu te dé el resultado exacto, sino de que no te dé un resultado bien pendejo.”}
\end{minipage}}
\end{center}

\end{document}